\documentclass[10pt,a4paper,twoside,bg=print,twocolumn,openany,nodeprecatedcode]{dndbook}
\usepackage[utf8]{inputenc}
\usepackage{tabulary}
\usepackage[normalem]{ulem}
\usepackage{color,soul}
\usepackage{float}
\usepackage{caption}
\usepackage{wrapfig}
\usepackage{tocloft}
\usepackage[hidelinks]{hyperref}

\raggedbottom
\extrafloats{2000}
\cftsetindents{section}{1em}{0em}
\cftsetindents{subsection}{2em}{0em}
\cftsetindents{subsubsection}{3em}{0em}
\setcounter{tocdepth}{2}
\definecolor {mytable} {HTML} {f4edd8}
\definecolor {mycomment} {HTML} {d0cfc2}
\definecolor {mysidebar} {HTML} {d0cfc2}
\definecolor {myreadaloud} {HTML} {d9e8f1}
\definecolor {myreadaloudborder} {HTML} {5a6173}
\definecolor {mypagenum} {HTML} {2a2d2e}
\DndSetComplexThemeColor[mytable][mycomment][mysidebar][myreadaloud][myreadaloudborder][titlered][titlegold][pagegold]

\captionsetup[table]{labelformat=empty,font={sf,sc,bf,},skip=0pt}

\begin{document}
\frontmatter

\tableofcontents
\newpage
\mainmatter%


\addcontentsline{toc}{chapter}{Introduction: Answering the Call}
\markboth{INTRODUCTION: ANSWERING THE CALL}{INTRODUCTION: ANSWERING THE CALL}
\chapter*{Introduction: Answering the Call}

\DndDropCapLine{L}{ong ago}, a war between gods of good and evil scoured the surface of the world of Exandria. Both the noble Prime Deities and the selfish Betrayer Gods imbued their mortal champions with supernatural gifts that granted them strength enough to challenge their divine foes. The greatest of these champions was Alyxian the Apotheon, who received the blessings of three Prime Deities to save the world from the apocalyptic power of the Betrayer Gods.

Centuries later, all knowledge of this godlike hero has been lost to time, and two groups of adventurers are called upon to retrace this mythic hero's steps and free him from his timeless prison. The tale begins in lands familiar to Critical Role fans—the war-scarred wastes of Xhorhas and the glimmering oasis city of Ank'Harel—but ultimately stretches beyond the Material Plane to a realm of lightless despair known as the Netherdeep.

\section{World of Exandria}
This adventure is set in Exandria, the world of Critical Role. Its continents include Tal'Dorei, Wildemount, Marquet, and Issylra. The adventure begins in a region of Wildemount called Xhorhas and transitions to the great Marquesian oasis city of Ank'Harel. More information about Exandria can be found in Explorer's Guide to Wildemount, but you don't need that resource to run this adventure.
The following sections describe events, regions, and other elements of Exandria that are significant to this adventure.

\begin{DndSidebar}[float=!b]{What Is Critical Role?}
Critical Role is a live-streamed show that airs on Thursday nights and stars voice actors Travis Willingham, Marisha Ray, Sam Riegel, Taliesin Jaffe, Ashley Johnson, Liam O'Brien, Laura Bailey, and Matthew Mercer. Each season is a complete Dungeons \& Dragons campaign that revolves around a particular band of lovably flawed heroes, with Matthew Mercer guiding the narrative and breathing life into the world of Exandria as the show's Dungeon Master. Like any D\&D game, the show is full of drama, laughter, silly voices, dice rolls, and stories that will live on in infamy.
\end{DndSidebar}

\subsection{The Calamity}
A cataclysmic war among gods and mortals occurred on Exandria centuries ago. This war, known as the Calamity, ended when the Prime Deities locked their evil brethren—and themselves—behind a Divine Gate that prevents the gods from directly interfering with the affairs of mortals. The gods can still bestow power upon their faithful and send powerful supernatural beings to help enact their will, but they can't directly touch the mortal works of Exandria.
Certain events of the Calamity are of paramount importance in this adventure, particularly the story of one forgotten hero who was favored by three of the Prime Deities and stood against the evil Betrayer Gods. As with other Critical Role stories, this tale involves a Vestige of Divergence, a magic item born out of the Calamity that grows in power alongside its bearer. This Vestige, the Jewel of Three Prayers, is described in the "Story Overview" section, and its magical properties are detailed in appendix B.
\subsection{Xhorhas}
Xhorhas, on the continent of Wildemount, was once a land dominated by the Betrayer Gods. It was their seat of power during the Calamity, and after their defeat, it was left a war-scarred wasteland prowled by demons. Today, it is a land populated by diverse creatures and people often considered to be monsters by their neighbor, the totalitarian Dwendalian Empire. The dominant political power in Xhorhas is the Kryn Dynasty, founded by drow who fled the Underdark and escaped the fell influence of Lolth the Spider Queen. They have discovered the power of an unusual deity of light and rebirth called the Luxon. The Bright Queen Leylas Kryn leads her people to glory in this blighted land.
Many beings lived in the wastes of Xhorhas before the drow of the Kryn Dynasty escaped the Underdark and built their civilization on the surface. Goblins, hobgoblins, bugbears, orcs, lizardfolk, kobolds, and countless other beings have banded together under the banner of the Kryn Dynasty—some willingly, others under duress.
Xhorhasian locations the characters will visit in this adventure include the settlements of Jigow and Bazzoxan, both described below.
\subsubsection{Jigow}
This coastal town in northern Xhorhas is populated mostly by goblins and orcs. A number of drow, many of them politicians and soldiers of the Kryn Dynasty, reside here as well. Tension exists between the people who value the ways of the village elders and those who prioritize the laws of the Kryn Dynasty. The town has a strong cultural tradition of competition, and the local Festival of Merit is the backdrop for chapter 1 of this adventure.
\subsubsection{Bazzoxan}
The military encampment of Bazzoxan keeps watch over the Betrayers' Rise, a dark temple that was instrumental to the Betrayer Gods during the Calamity. Xhorhasian explorers and settlers built the foundation of a new city around the looming stone structures of the abandoned sanctum—until demonic forces from within the old temple threw the burgeoning society into chaos.
As the characters discover in chapter 3, Bazzoxan is now locked in a stalemate between the Kryn Dynasty's forces and the demons of the Betrayers' Rise. This unholy site is the subject of constant research and of growing worry for the dynasty.
\subsection{Ank'Harel}
The continent of Marquet is scarred by a vast, arid desert, which was created long ago when the Betrayer God known as Gruumsh the Ruiner struck a cataclysmic blow with his spear into Marquet's lush jungles. Within this desert is one of the greatest city-states known to the mortals of this age: Ank'Harel, also called the Jewel of Hope.
Traders who needed to traverse the desert to sell their wares founded Ank'Harel around an oasis, unknowingly locating their new home above the sunken ruins of an elven city called Cael Morrow, which was destroyed by Gruumsh in the same spear thrust that laid waste to much of the continent.
Several factions are active in Ank'Harel, and the characters can receive help on their mission by joining one of these groups. Descriptions of the factions are presented in chapter 4.
\subsection{Cael Morrow}
Beneath Ank'Harel lie the remains of a former Marquesian metropolis. Its true name is lost to time, since nearly all records from before the Calamity were obliterated when Gruumsh destroyed the city. Modern scholars call it Cael Morrow, the Drowned City. The epithet refers to the fact that the ruins are submerged in a vast underground cistern—part of the same body of water that feeds the oasis that gives life to the city.
Hidden within Cael Morrow is a rift to the Netherdeep, which the characters discover in chapter 5.
\subsection{The Netherdeep}
A realm of despair, the Netherdeep was created when the spear of Gruumsh pierced through the Material Plane and into the Elemental Plane of Water. It is an extraplanar labyrinth of lightless trenches, infused with the distorted emotions of a demigod who is both the Netherdeep's prisoner and jailer. This location awaits the characters in chapter 6.

% \RequirePackage{calc}

% \newlength{\titleborderwidth}

% %Set border width
% \setlength{\titleborderwidth}{0.16\paperwidth}

% %Set page geometry
% \geometry{inner=\titleborderwidth}

\begin{tikzpicture}
    \node[anchor=north east, align=left, inner sep=0] at (current page.north east){
        \includegraphics[width=\linewidth]{images/adventure/CRCotN/000-00-001.ruidus-moon.png}};
\end{tikzpicture}
\subsection{Ruidus}

% \includegraphics[width=0.5\textwidth]{images/adventure/CRCotN/000-00-001.ruidus-moon.png}

Two moons hang in the sky of Exandria. One is Catha, a large, pale orb that cycles through its phases once per month. The other is Ruidus, a ghostly vermilion satellite that circles Exandria once every six months.
Catha's brightness is constant, its phases predictable. Conversely, Ruidus exhibits strange behavior, seeming to glow more brightly at times or suddenly casting off its shadow to appear full. The vermilion moon often disappears from the night sky entirely, while at other times it unexpectedly flares with bright, ruddy light.
Though this story presents Ruidus as a subject of superstition and folklore, it intentionally avoids discussing the moon's actual magical power. Any mystical powers that nonplayer characters in this story claim to receive from the vermilion moon are related to Ruidus only so far as people believe them to be. In truth, the powers those characters wield comes from a figure called the Apotheon. The true nature of Ruidus is a topic to be explored in other Critical Role stories.
Not all people fear Ruidus, but superstitions about it are widespread across Exandria. Some of the more well-known superstitions include the following:

\begin{description}[nolistsep]
\item[Harbinger.]
Those who are born on a night when Ruidus is full are destined to bring suffering to others or to experience great tragedy in their own lives.
\item[Malice.]
Those who study Ruidus and become obsessed with its secrets are compelled to cause misfortune and woe.
\item[Misadventure.]
Plans made or set in motion when Ruidus is full are doomed to failure, often due to betrayal or miscommunication.
\end{description}

\begin{tikzpicture}
    \node[anchor=north west, align=left, inner sep=0] at (current page.north east){
        \includegraphics[width=\linewidth]{images/adventure/CRCotN/001-00-002.alyxian.png}};
\end{tikzpicture}

% \includegraphics[width=0.5\textwidth]{images/adventure/CRCotN/001-00-002.alyxian.png}

\section{Story Overview}
The story at the heart of this adventure begins long ago in an age shrouded in myth and legend, as befits an epic Critical Role tale.


\subsection{Rise and Fall of the Apotheon}
During the Founding, a time when the gods still walked the face of Exandria, the world's divine creators discovered an unidentifiable power seeping through the fabric of reality. Legends assert that this alien influence was a threat to all life on Exandria, and the gods banded together to banish it.
This cancerous incursion of dark power is said to have crystallized into Ruidus, the small, vermilion moon that hangs in the sky along with Catha, the world's natural moon. The gods agreed to create a tale about Ruidus to conceal its alien origin from the mortals of the world, informing them that it was a moon of ill omen, and its magical influence was always to be avoided. This tale concocted by the gods was not a lie, for Ruidus's alien magic twists the fate of those who are born or embark on ventures while bathed in its vermilion light.
Some eons later, a schism rocked the pantheon of Exandria, splitting the divine beings into Betrayer Gods and Prime Deities. The Betrayer Gods strove to ruin and dominate the world, and they battled the Prime Deities in a war called the Calamity.
Near the beginning of this war, a boy named Alyxian was born in the lands of Wildemount, the seat of the Betrayer Gods' power, beneath the full light of Ruidus. He was said to be cursed from birth and was viewed with suspicion by all around him, except for his parents. As a young man, he sought to defy his fate and distinguish himself on the field of battle. Alyxian found himself embroiled in the most vicious battles of the war. He fought selflessly and begged the Prime Deities to help him protect the innocent people caught up in their war. Thrice he asked, and thrice they answered, granting him greater and greater power. He was more than a champion of the gods; he had become the Apotheon—a noble warrior so suffused with divine energy that he was halfway to godhood himself. In those times, he wore a gods-gifted amulet known as the Jewel of Three Prayers as a symbol of his divine covenants.
Eventually, the war brought the Apotheon to the great jungle city now known as Cael Morrow, a utopian domain of elves and orcs. These orcs were the first generation of their kind, who had been transformed by the spilled blood of Gruumsh the Ruiner. Although fear had reigned for a short time after the transformation, the two peoples soon realized they were still kin. The city was the jewel of the lush and verdant continent of Marquet, and its people gave thanks daily to Corellon the Arch Heart for their protection.
The Ruiner despised the Arch Heart and swore to annihilate Cael Morrow and all of Marquet with a single stroke of his spear. He strode into the center of the city, smiting all who dared raise a blade against him. When he reached the city's heart, Gruumsh raised his spear to strike the earth with cataclysmic force—but his blow was intercepted by the Apotheon at the last possible second.
The Apotheon's intervention kept Gruumsh from annihilating all life on Marquet, but the force of his spear thrust still brought the lands of Marquet to ruin. Fire raged across the lush jungle, turning verdant beauty into blasted desert, and Cael Morrow was shunted deep into the earth. Towers toppled, stone crumbled, and all trace of the great civilization was wiped from the face of Exandria.
In that moment of destruction and death, the Apotheon's connection to Ruidus flared to life. A rift was torn between worlds, in which alien energy from an unknown realm and the waters of Cael Morrow's oasis mingled to produce a lightless realm of water and strange magic. There, in what came to be known as the Netherdeep, the Apotheon has been trapped for untold ages, consumed with sadness, furious over his defeat, and yearning for his freedom.
In time, the immortal Apotheon fell into a long and troubled slumber. In his dreams, the barren caverns of the Netherdeep began to shift. His memories filled the darkness, and a cocoon of melancholy formed around his body. From this Heart of Despair emerged a crimson element that embodied the Apotheon's power. It spread, growing and crystallizing as it moved, until all of the Netherdeep was suffused with its power.
\subsection{The Apotheon Awakens}
Centuries have passed since the Calamity, and life has returned to the desert lands of Marquet. A grand city called Ank'Harel has been built around a desert oasis. Little did its founders know that the water of their oasis originated in the Netherdeep, and that beneath the new desert metropolis was an underwater cavern holding the ruined city of Cael Morrow.
A new form of conflict entered the story when the Allegiance of Allsight, an influential group of academics in Ank'Harel, discovered a strange mineral in the sunken ruins beneath their city. This substance—a slick, oily stone veined with blood-like streaks—possessed unknown magical properties. The Allegiance tried to keep its discovery secret, but the news soon attracted the attention of a rival faction.
The Consortium of the Vermilion Dream—a secret occult society obsessed with the magic of the Moon of Ill Omen, Ruidus—heard a rumor about the discovery of this mineral. Consortium agents broke into the Allegiance of Allsight's excavation site in Cael Morrow and learned about the discovery firsthand. They dubbed the mineral "ruidium," because its red veins and its mysteriousness reminded them of the light and power of Ruidus.
The consortium continued sneaking into the excavation, hunting further deposits of the mineral, and discovered a planar rift within the ruins. One of their agents succeeded in getting inside the rift—and when that agent returned, she claimed that she had entered a dark, underwater realm where the walls themselves were veined with ruidium. She managed to recover a sample and escape before the water's intense pressure crushed her body.
When that sample of ruidium was torn from the Netherdeep, the shock of the extraction roused the dormant Apotheon from his slumber. All at once, centuries of accumulated grief crashed down upon him, and the Netherdeep roared with the strength of his tempestuous emotions.
The power that flows from the Apotheon is no longer truly his own. Try as he might to hang on to the memory of the hero he once was, he is driven entirely by his tortured emotions. Throughout his centuries of isolation, the dormant, dreaming Alyxian had hoped for someone to find and rescue him. But when mortals discovered him, they seemed interested only in stealing his power for themselves. The furious Apotheon tapped into the alien nature of the Netherdeep to cause ruidium to slowly corrupt all who hold it.
Alyxian beseeched the three gods he had prayed to in ancient times—Sehanine the Moon Weaver, Avandra the Change Bringer, and Corellon the Arch Heart—to send heroes to save him. Despite being sealed behind the Divine Gate, the gods were able to propel the Jewel of Three Prayers to a site where Alyxian once beseeched the Moon Weaver for aid long ago. The jewel materialized within a sunken temple in the Emerald Gulch near Jigow, and the gods hoped that a good-hearted band of heroes would find it, learn of the Apotheon's plight, and rescue their tortured champion. The Jewel of Three Prayers is currently dormant, waiting to be awakened by a determined hero.
\subsection{The Apotheon's Call}
The first chapter of this adventure begins in Jigow, a Xhorhasian town populated mostly by orcs and goblins brought together by the Kryn Dynasty. Smaller populations of other peoples live among them. Unlike some other Kryn settlements, the people of Jigow have maintained their city's distinctive culture. Jigow's most popular event is the annual Festival of Merit, in which locals and visitors compete in tests of strength and cunning. Rivalries are forged, bets are won and lost, and bragging rights are secured for the year—until the next festival.
During this event, destiny calls to two adventuring parties: the characters (run by your players) and a group of rival adventurers (run by you). The grand finale of the Festival of Merit sends both groups into an underwater grotto. Here, the characters discover a place where Alyxian prayed to Sehanine the Moon Weaver centuries ago. Here, the Jewel of Three Prayers waits to be found. The discovery of this item and a vision in which Alyxian pleads with the characters to save him from his prison of darkness and despair set the adventure in motion. If the characters answer the call, they will come into conflict with the power-hungry factions that awakened the Apotheon, the alien evil of the Netherdeep itself, and the rivals who have dogged their path since the festival at Jigow.
\section{Running the Adventure}
Call of the Netherdeep is a Dungeons \& Dragons adventure optimized for five characters. The player characters are the heroes of the story; this book describes the villains and monsters the heroes must overcome and the locations they must explore to bring the adventure to a successful conclusion.
To run the adventure, you need the fifth edition core rulebooks (Player's Handbook, Dungeon Master's Guide, and Monster Manual).
\begin{DndReadAloud}{}
Text that appears in a box like this is meant to be read aloud or paraphrased for the players when their characters first arrive at a location or under a specific circumstance, as described in the text.
\end{DndReadAloud}

When a creature's name appears in bold type, that's a visual cue pointing you to its stat block as a way of saying, "Hey, DM, you better get this creature's stat block ready. You're going to need it." Usually, you can find the stat block in the Monster Manual; if the stat block appears elsewhere, the text tells you so.
Spells and equipment mentioned in the adventure are described in the Player's Handbook. Magic items are described in the Dungeon Master's Guide unless the adventure's text directs you to an item's description in appendix B of this adventure.
\subsection{Character Advancement}
The characters begin the adventure at 3rd level. If your players are creating 3rd-level characters from scratch, assume that each character has the normal starting equipment for their background and class.
If you want to run this adventure as part of a campaign that starts at 1st level, the adventure "Unwelcome Spirits" in Explorer's Guide to Wildemount is ideal for a group of 1st-level characters in Xhorhas and can be used to advance the characters to 3rd level.
The Character Level Advancement table summarizes when the characters gain levels during this adventure.
\begin{DndTable}[header=Character Level Advancement]{XX}

Level & Requirement\\
4th & Finish the Emerald Grotto challenge in chapter 1.\\
5th & Arrive at Bazzoxan at the end of chapter 2.\\
6th & Defeat or frighten away the gloomstalkers that escape from the Betrayers' Rise in chapter 3.\\
7th & Reach the prayer site in the Betrayers' Rise (area R16) at the end of chapter 3.\\
8th & Complete at least three faction missions in chapter 4.\\
9th & Complete at least one faction mission in chapter 4 that requires the party to visit the sunken ruins of Cael Morrow.\\
10th & Enable the Jewel of Three Prayers to transform into its Jewel of Three Prayers (Exalted) (see area M9).\\
11th & Obtain at least three Fragments of Suffering by exploring the Netherdeep in chapter 6.\\
12th & Enter the Heart of Despair at the end of chapter 6.\\
13th & Defeat Alyxian or enable him to be redeemed in chapter 7.\\
\end{DndTable}


\subsection{Adventure Structure}
Call of the Netherdeep is divided into seven chapters, as shown in the Adventure Flowchart.
Flowchart not implemented -- yet.
\section{Roleplaying the Apotheon}
Alyxian is trapped in an extradimensional prison of his own making and longs for freedom. He still has a connection to several holy sites where he received power from the gods during the Calamity, and his consciousness reaches out to certain mortals whose hearts resonate with his own.
Beyond this, Alyxian yearns to be remembered. He is a being composed of little more than memory, sorrow, and hatred. He fears that the world has forgotten his sacrifice entirely—and the Netherdeep has preyed on his selfish urges, instilling in him an obsessive desire to be remembered and idolized by a world that seemingly has forsaken him.
Alyxian appears to the characters several times in visions. The version of himself he shows them is the idealized way he sees himself—someone who has always tried to do good but has been punished nonetheless. The Apotheon is trying to put his best self forward to get what he wants out of the characters. This isn't a false persona, but it isn't the whole picture. In later visions, when the characters are able to speak to Alyxian and ask him questions, probing too deeply can cause the Apotheon to reveal the more raw, selfish parts of himself. These parts of him see the characters as pawns—tools granted to him by the gods.
Ultimately, your job is to try to balance the selfless man that Alyxian once was with the traumatized, manipulative being that he has become. By the end of the adventure, if all goes well, the characters will realize that Alyxian can be redeemed, and that helping him overcome his pain and trauma—thus restoring the hero who once saved Exandria from destruction—is a cause worth fighting for.
\section{Story Tone and Player Limits}
Like other Critical Role stories, this adventure walks a line between optimistic heroism and morally ambiguous character dilemmas. Heroic and equivocal choices are contrasted with moments of supernatural evil. Although this adventure is not a horror tale, it does involve elements of fear, suspense, and the grotesque.
The events of the adventure should make the characters feel stressed and anxious, but the players should be relaxed and having fun. Before starting the adventure, have a candid conversation with your players about hard and soft limits on what topics can be broached at the table. Your players might have phobias and triggers you aren't aware of; never assume you know your players' deepest fears, no matter how long you've been playing together. Any topic or theme that makes a player feel unsafe should be avoided. If a topic or theme makes one or more players nervous but they give you consent to include it in-game, incorporate it with care. Since D\&D adventures aren't heavily scripted, these elements might arise unexpectedly during play. Be prepared to move away from such topics and themes quickly if any player feels uncomfortable.
\section{Underwater Adventuring}
During this adventure, the characters will explore sunken grottoes, submerged ruins, and an extraplanar aquatic abyss. To run these underwater scenarios in a smooth and fun manner, familiarize yourself with the underwater combat rules in the Player's Handbook as well as the rules below.
\subsection{Swimming}
Unless aided by magic, a character without a swimming speed can't swim indefinitely. After each hour of swimming, a character must succeed on a DC 10 Constitution saving throw or gain 1 level of exhaustion.
A creature with a swimming speed—including a character with a ring of swimming or similar magic—can swim all day without penalty and uses the normal forced march rules in the Player's Handbook.
Swimming through deep water is similar to traveling at high altitudes because of the water's pressure and cold temperature. For a creature without a swimming speed, each hour spent swimming at a depth greater than 100 feet counts as 2 hours for the purpose of determining exhaustion, and swimming for an hour at a depth greater than 200 feet counts as 4 hours.
\subsection{Underwater Visibility}
Visibility underwater depends on water clarity and the available light. Unless the characters have light sources, use the Underwater Encounter Distance table to determine the distance at which characters underwater become aware of a possible encounter.
\begin{DndTable}[header=Underwater Encounter Distance]{XX}

Water Clarity and Light & Encounter Distance\\
Clear water, bright light & 60 ft.\\
Clear water, Vision and Light & 30 ft.\\
Murky water or Vision and Light & 10 ft.\\
\end{DndTable}


\subsection{Netherdeep Water Pressure}
The Netherdeep has the added danger of extreme water pressure, the effects of which are described in chapter 6.
\section{Ruidium}
The Apotheon's distorted emotions manifest in the form of a crimson element called ruidium, so named because its color is similar to that of Ruidus.
Ruidium originates in the Netherdeep. It penetrates and grows into all things it touches. Inanimate objects that become bonded with ruidium are imbued with a sliver of the Apotheon's power. Creatures that interact with ruidium can be physically and emotionally corrupted by it. Their bodies become laced with ruidium veins, and their minds become tormented by the same emotions that haunt the Apotheon.
\subsection{Ruidium Corruption}
When a creature is at risk of becoming corrupted, it must make a Charisma saving throw. The save DC varies based on the ruidium's virulence. On a failed saving throw, the creature gains 1 level of exhaustion (see the Player's Handbook) and becomes corrupted if it isn't already. From that point on, the creature takes 1d10 psychic damage whenever its level of exhaustion increases or decreases by 1, until it's no longer corrupted (see "Ending Corruption" below).
\subsubsection{Signs of Corruption}
The first physical sign of corruption in a creature is a bright red rash, which appears on the creature's body where it made contact with the ruidium. As the creature's level of exhaustion increases, the signs of corruption become more obvious, as summarized in the Physical Signs of Ruidium Corruption table. Decreasing a creature's level of exhaustion doesn't affect the physical signs of corruption. Once a new physical sign appears, it can't be removed from the creature until its ruidium corruption is ended.
\begin{DndTable}[header=Physical Signs of Ruidium Corruption]{XX}

Exhaustion Level & Physical Signs\\
1 & A red rash appears, originating from the point of contact with ruidium.\\
2 & Pulsing crimson veins spread across the creature's skin.\\
3 & Crimson blisters and boils appear on the creature's skin.\\
4 & Stubby spurs of ruidium crystal protrude from the creature's body.\\
5 & The creature's crystal protrusions grow larger and more grotesque.\\
6 & The corruption kills the creature.\\
\end{DndTable}


A creature beset by ruidium corruption also exhibits emotional signs of corruption that worsen as its level of exhaustion increases. These symptoms include amplified feelings of regret, yearning, rage, and despair. A player whose character is corrupted by ruidium can roleplay these symptoms however they wish. (For example, the player could emphasize or amplify their character's flaw, or choose a new flaw for their character.) As with the physical signs of corruption, all emotional symptoms persist until the creature's ruidium corruption is ended.
\begin{DndSidebar}[float=!b]{Ruidium Spell Components}
One ounce of ruidium can be substituted for 500 gp worth of any material component needed to cast a spell. A creature that casts a spell using ruidium as a replacement component must succeed on a DC 20 Charisma saving throw or gain 1 level of exhaustion after the spell is cast. If the creature isn't already suffering from ruidium corruption, it becomes corrupted if it fails the saving throw.


\end{DndSidebar}

\subsubsection{Ending Corruption}
Only a wish spell or divine intervention can cure a creature's ruidium corruption; simply removing all levels of exhaustion from the creature cannot.
Ruidium and the effects of its corruption vanish from the world if the Apotheon is redeemed at the end of this adventure (see "Best Ending: A World that Remembered" in chapter 7).
When ruidium corruption ends on a creature, all physical and emotional signs of it disappear from the creature instantly.
\section{Rivals}
From the beginning, the characters are challenged by a rival adventuring party. This rivalry starts small, with the two groups first meeting while they play games of skill and strength in Jigow. As the rivals make their own way through the adventure, they grow in power and ambition just as the characters do. They are determined to prove that they are the true heroes of this story.
This adventure presents five rivals. Ideally, the number of rivals should be equal to the number of characters in the players' party. If there are fewer than five characters in their party, remove one rival of your choice or give one of them a dramatic death early on. If there are more than five characters in your party, create new rivals. (Maybe one of the rivals has an identical twin.) Alternatively, start with a party of five rivals, and keep an eye out for chances to add new nonplayer characters to the rival party as the characters meet them.
The following descriptions pertain to the rivals at the start of the adventure.
\includegraphics[width=0.5\textwidth]{images/adventure/CRCotN/002-00-003.ayo-jabe.png}
\subsection{Ayo Jabe}
As the de facto leader of a new, unnamed adventuring party, Ayo Jabe has a lot of responsibility on her shoulders. She and her companions have worked odd jobs around the town of Jigow for a few weeks now, and she's starting to feel confident that they're ready for real adventure.
Ayo is a water genasi who was born to orc parents during a terrible storm in the Emerald Gulch. Her family has lived in Jigow for generations, and she's eager to leave home and see the world.
Ayo picked up her combat skills while working as a hunter for the town of Jigow. She recently joined up with her childhood friend Dermot, as well as Maggie Keeneyes, a mercenary who came to town several weeks ago. Despite Maggie's brusque demeanor, she became a loyal friend of Ayo and Dermot, and the three of them have become inseparable, even as more adventurers have joined their party.
Ayo is hotheaded and appreciates people who make decisions as impulsively as she does. Nevertheless, she respects Dermot's even-tempered advice and dutifully plays the role of the wise leader. She has no patience for people who are indecisive, and she is infuriated by people who don't say what they really mean.
\includegraphics[width=0.5\textwidth]{images/adventure/CRCotN/003-00-004.dermot-wurder.png}
\subsection{Dermot Wurder}
When you're the twelfth son of a poor goblin family in Jigow, the only way to make a name for yourself is to become a great champion—someone who can win bragging contests in the local taverns night after night. Young Dermot Wurder, however, wasn't interested in performing feats of strength or agility that would win him a boast-worthy epithet. Fame wasn't for him, nor was the aggressiveness that becoming famous required. He was more interested in cooking, planting flowers, and sewing clothes—doing the work that kept his family together while his carefree siblings dove off waterfalls and wrestled stray dogs.
Dermot and Ayo Jabe have known each other long enough that their first meeting has vanished into the haze of youthful pre-memory. Dermot accompanied Ayo whenever he could, packing herbs and medicines when they went exploring in the woods. When Ayo recently told him of her dream to form an adventuring party, he started training to wear heavy armor and threw himself into studying the faith of the Luxon so he could wield the blessings of the light to protect those he cares about.
Dermot is fiercely protective of his friends and furiously rebukes anyone who disparages or threatens them. His deepest need—one even he doesn't know he has—is to make a friend who will help him realize what he wants for himself.
\includegraphics[width=0.5\textwidth]{images/adventure/CRCotN/004-00-005.galsariad-ardyth.png}
\subsection{Galsariad Ardyth}
Beautiful, ethereal, deathly, shadowy—all accurately describe Galsariad Ardyth, a drow in his two-hundredth year of life. He's recently taken up the study of arcane magic, and he's pursuing the life of an adventurer in hopes of improving his reputation within the Kryn Dynasty. Loquacious, snarky, and sarcastic to a fault, he's ready with a barb for any occasion—usually to mask his own insecurities.
A city-dweller from the Kryn capital of Rosohna, Galsariad is blessed with sharp aesthetic sensibilities and an interest in ancient lore, especially history concerning the Age of Arcanum, Exandria's long-lost magical golden age. He's also the newest member of the rival party, and both Ayo and Irvan are impressed by his ethereal elegance—and have a bit of a crush on him, even if he does talk too much.
Galsariad appreciates people who share his interests and are willing to spend time studying with him, though he also likes it when people treat him with respect even if they don't care about magic. He dislikes people who keep secrets from him, and hates when people judge him for being a novice at magic even though he's centuries old.
\includegraphics[width=0.5\textwidth]{images/adventure/CRCotN/005-00-006.irvan-wastewalker.png}
\subsection{Irvan Wastewalker}
The name "Wastewalker" conjures fear on the plains of Xhorhas, as the name of a clan of Xhorhasian nomads. Irvan—Irv to his friends—was born into that clan, but he hates the name. He left his family when he was a teenager, struck out on his own across the wastes, and made his way to the city of Asarius. Even though the "City of Beasts" is renowned for danger, Irv felt at home there in a way he'd never felt with his clan.
He earned a reputation for being a party animal who could put down a dozen drinks in a night and still dance with the vigor and elegance of an archfey. Irv is young—just a few weeks past his nineteenth birthday—and he is secretly ashamed of his patchy beard and scrawny, human body. He has a bigger secret related to that feeling: he was consecuted in another life. The act of consecution is a sacred Kryn ritual of rebirth through the mystical power of the Luxon. Irvan has memories of his previous life as a bugbear of Den Hythenos, and he hopes to prove to himself that he can do great things without the aid of his powerful "family" in the Kryn Dynasty.
He met Ayo and her friends when he traveled to Jigow to experience its contests of strength and skill. He sticks with the group in part because he's interested in what adventures they'll get into, and also because he's smitten with the group's newest addition, a drow arcanist named Galsariad.
\includegraphics[width=0.5\textwidth]{images/adventure/CRCotN/006-00-007.maggie-keeneyes.png}
\subsection{Maggie Keeneyes}
People might laugh when a 12-foot-tall ogre orders a drink at a bar and says her name is Maggie, but they don't laugh for long. Some people fixate on her name, her enormous size, her muscles, or the weapon at her side. Wiser folk take notice of Maggie's bright blue eyes. All her life, people have considered Maggie a stupid meathead because of her size, but her eyes betray her intelligence. She can read others with a glance, whether in conversation or in a duel. When her eyes dart back and forth across a battlefield, they take in enough information to give her allies an advantage in the fight.
Things changed for Maggie when she first arrived in Jigow and met Ayo Jabe three weeks ago. She came looking for mercenary work to make ends meet but found a true friend instead. Ayo saw the intelligence in Maggie's eyes and was keen to hear Maggie's thoughts. They became fast friends, and Maggie would sooner die than let harm come to her new companion.
Maggie loves poetry and music with profound lyrics, as well as matching wits with people as clever as she is. She hates being stereotyped and has a dim view of those who are too quick to judge others.
\subsection{Rivals' Goals}
Each rival has goals. As the characters interact with the rivals, they might choose to help the rivals achieve their goals or use their own knowledge of those goals to manipulate the rivals into doing what they want. A Charisma check to ask for help (see "Asking for Help" below) is made with advantage if the rival thinks agreeing to the character's request will help achieve the rivals' goals and is made with disadvantage if it goes directly against the rivals' goals.
The Rivals' Goals table summarizes each rival's goals at the start of this adventure. Some of them are broadly applicable and might never be definitively achieved. Others are highly specific (such as Galsariad's goal to match wits with an archmage). The rivals can acquire new goals as the adventure progresses, as you see fit.
A character can learn about one of a rival's goals by spending at least 1 minute in conversation with that rival and making a DC 10 Wisdom (Insight) check. No check is necessary if the rival is friendly toward the character and divulges the information willingly.
\begin{DndTable}[header=Rivals' Goals]{XX}

Rival & Goals\\
Ayo Jabe & Become a hero like the ones she has read about in stories; make a friend she can be truly herself around; kill a legendary monster\\
Dermot Wurder & Protect his friends; have a life-changing holy vision; discover that he has value himself, not just as someone who can help others\\
Galsariad Ardyth & Match wits with an archmage; discover a magical secret no one else knows; be respected by someone he respects\\
Irvan Wastewalker & Experience things he has never encountered before; fall in love with someone who doesn't know about his past lifetime; appear mature despite his youthful appearance\\
Maggie Keeneyes & Spar with a true tactical genius; write a song or poem that causes someone to weep with emotion; be able to retire and never kill again\\
\end{DndTable}


\subsection{Rival Stat Blocks}
Stat blocks for the rivals can be found in appendix A. Each of the rivals has three stat blocks; like the characters, they become more powerful as the adventure progresses. The Rival Stat Blocks table shows you which stat blocks to use based on the chapter you are running.
\begin{DndTable}[header=Rival Stat Blocks]{XX}

Chapters & Stat Blocks\\
1 and 2 & Tier 1\\
3 and 4 & Tier 2\\
5, 6, and 7 & Tier 3\\
\end{DndTable}


\subsection{Relationships with the Rivals}
When the characters meet the rivals in chapter 1, all the rivals have an indifferent attitude toward them. It's up to you to determine when a rival's attitude toward a character might change. For example, after a combat encounter in which one character saves a rival's life or helps a rival achieve one of their goals (see the Rivals' Goals table), you can decide if that deed enabled the rival to see the character in a new light, changing the rival's attitude toward the character from hostile to indifferent or indifferent to friendly. The rules for changing attitude are detailed in the Dungeon Master's Guide.
\subsubsection{Keeping Track of Attitudes}
Since every rival's attitude toward the characters starts out as indifferent, it's important to note only the hostile and friendly relationships that blossom between the two parties and the reasons why the relationship became friendly or hostile.
You can decide that certain actions—such as a character killing one of the rivals—automatically turn all the rivals hostile against all the characters. The converse isn't true; singular acts of kindness toward an individual rival aren't as likely to make all the rivals friendly toward all the characters.
\begin{DndSidebar}[float=!b]{Vital Nonplayer Characters}
Critical Role stories are known for their rich cast of characters. The characters created by the players are the protagonists of this story. No matter how important this adventure's demigod, power-hungry faction leaders, and precocious rival adventurers might seem to be, it's the player characters' story first and foremost.


Nevertheless, these other participants are of crucial importance. They exist to goad the heroes into making reckless decisions, to prey on their flaws, to spur the story forward when it stalls, and to enhance the mysteries of the world.


\end{DndSidebar}

\subsubsection{Asking for Help}
At times, the characters might want help from one or more rivals to accomplish their goals. Any character who requests help from a rival must make a successful Charisma check to convince that rival to help them. Consider that rival's relationship to the character in question, then consult the Conversation Reaction table in the Dungeon Master's Guide to determine the DC of the check and the rival's reaction to the character's request.
\subsubsection{Coming to Blows}
As long as the characters haven't slain any of them, the rivals do their best to avoid killing the characters. If a battle erupts between the two groups, the rivals try to knock the characters unconscious (see the rules for doing so in the Player's Handbook). In situations where that's not possible, Dermot Wurder uses his spare the dying spell to stabilize a dying character.
If the heroes surrender, Ayo is usually happy to accept that outcome. She would prefer to let her foes go free, though she demands they stay out of her group's way.
You determine if the situation is dire enough for the rivals to kill the characters or take them prisoner, or if they would even refuse to accept a surrender.

\begin{figure*}[t]
\includegraphics[width=\linewidth]{images/adventure/CRCotN/007-00-008.xhorhas-adventurers.png}
\end{figure*}



\section{Character Creation}
This adventure begins in Xhorhas, where most humans are nomads; most elves are drow who live aboveground; dwarves and halflings are rare; and goblinoids, orcs, lizardfolk, kobolds, and other creatures sometimes seen as "monstrous" elsewhere are more populous than gnomes and dragonborn. Whether the characters are natives of Xhorhas or outsiders, assume they are aware of what's said in the "What the Characters Know" section below.
\subsection{Session Zero}
Before starting the adventure, consider spending a game session—colloquially called session zero—to establish expectations, share house rules, and help your players create characters. If you have a copy of Tasha's Cauldron of Everything, see that book's "Session Zero" section for advice on how to lay the groundwork in session zero for a fun campaign that exceeds your players' expectations.
All the character races presented in the Player's Handbook are well suited for this adventure. If your players have access to Explorer's Guide to Wildemount, the race and subclass options presented in that book are available to them as well. With your consent, a character can also receive the Hollow One supernatural gift described in that book.
Additional race options appear in Monsters of the Multiverse, with bugbear, duergar, goblin, hobgoblin, kobold, lizardfolk, minotaur, and orc being appropriate choices for players who want their characters to come from cultures or societies with roots in Xhorhas.
All class options presented in the Player's Handbook can be used in this adventure, and players can choose options from other D\&D books with your approval.
\begin{DndSidebar}[float=!b]{Making Mistakes}
Dungeon Masters are fallible, just like everyone else, and even the most experienced DMs make mistakes. If you overlook, forget, or misrepresent something, correct yourself and move on. This is a big adventure with lots of interconnecting elements; no one expects you to internalize or memorize every aspect of it. As long as your players are having fun, everything will be just fine.


\end{DndSidebar}

\subsubsection{Character Back Stories}
The beginning of this adventure assumes the characters are familiar with each other. Are they using Jigow as a rest stop on their way elsewhere, are they residents of Jigow, or some combination of the two?
If the players are unfamiliar with Exandria and the lands of Xhorhas, they might find it daunting to create a back story set in this world. If you have Explorer's Guide to Wildemount, the Heroic Chronicle section in that book roots a character's history in the lands of Wildemount by helping the player determine their character's birth nation, home settlement, and relationships with family members, allies, and enemies. It also helps establish major events that happened to the character before the campaign began.
Regardless of whether you use that information, asking the players the following questions about their characters can help them come up with a strong back story:

\begin{itemize}
\item What's one small heroic deed you've already accomplished?
\item Do you give your trust easily, or do people have to earn your trust?
\item What's one thing you want more than anything else? How will being an adventurer help you achieve it?
\end{itemize}

\subsubsection{What the Characters Know}
Characters begin the adventure knowing the following facts, which you can share with the players as they create their characters.
\paragraph{Kryn Dynasty}
The Kryn Dynasty is the dominant nation in Xhorhas. It was founded by a drow queen named Leylas Kryn, who fled the Underdark and the tyrannical rule of Lolth the Spider Queen along with her disciples. The Bright Queen still rules the dynasty centuries later, and its cities contain more than just drow. Orcs, goblinoids, tieflings, humans, and many others call the cities of the dynasty their home. Countless more denizens of the dynasty are nomads who roam the wastes in clans, hunting mastodons and other Xhorhasian megafauna.
\paragraph{Jigow}
This coastal settlement is actually a string of villages that are home to a collection of folk from all over Xhorhas. Goblin and orc clans founded Jigow, which explains why the settlement is governed by two elders—a goblin and an orc. The Aurora Watch (the military arm of the Kryn Dynasty) maintains a presence here, under the command of a drow called Taskhand Durth Mirimm.
Townsfolk tend to be competitive, and friendly rivalries are commonplace. Most of Jigow's residents live in a central region called the Jumble.
\paragraph{The Luxon}
The official deity of the Kryn Dynasty, whose symbol appears on the nation's heraldry, is the Luxon. This mysterious divine entity of light and rebirth has granted its faithful several esoteric secrets, the greatest of which is consecution—the act of preparing one's soul for rebirth. Through consecution, some people within the Kryn Dynasty have lived many lifetimes, often in bodies different from the ones they were first born in. In the sequence of consecution, a drow might become a goblin, then be reborn as a bugbear, then an orc, and so on—all the while gaining greater knowledge about the world through their experiences. This process has no mechanical benefit, but players can make consecution and rebirth an interesting part of their characters' back stories.
If a follower of the Luxon who has undergone the ritual of consecution dies within 100 miles of a Luxon beacon, their soul is ensnared by it and reincarnated within the body of a random Humanoid newborn within 100 miles of the beacon.
\section{Pronunciation Guide}
The Pronunciations table shows how to pronounce many of the unusual names that appear in this adventure.
\begin{DndTable}[header=Pronunciations]{XXX}

Name & Pronunciation & Description\\
Aloysia Telfan & ah-LOY-see-uh TELL-fin & Elf agent of the Consortium of the Vermilion Dream\\
Alyxian & ah-LICK-see-in & Tragic human hero who sacrificed himself to save others\\
Ank'Harel & AWNK-har-el & Majestic oasis city in Marquet\\
Apotheon & uh-PAW-thee-awn & Alyxian's title, denoting his godlike powers\\
Aradrine the Owl & AWR-uh-dreen & Goliath coleader of the Consortium of the Vermilion Dream\\
Ayo Jabe & AYE-oh JAW-bay & Water genasi leader of the rival party\\
Bautha Dyrr & BOUGH-thuh DEER & Drow priest in Bazzoxan\\
Bazzoxan & BAZ-oh-zan & Military outpost in Xhorhas\\
Beltreath & BELL-treeth & Sword wraith commander in Cael Morrow\\
Bookkeeper Khime & KEEM & Orc member of the Allegiance of Allsight\\
Cael Morrow & KAYL MAW-row & Sunken ruin of an ancient city beneath Ank'Harel\\
Catha & KATH-uh & Exandria's larger, silvery moon\\
Colbu Kaz & COAL-boo KAZ & Goblin elder of Jigow\\
Dendarron the Sun Bear & den-DAR-run & Halfling coleader of the Consortium of the Vermilion Dream\\
Dermot Wurder & DER-mott WER-der & Goblin priest of the Luxon and a member of the rival party\\
Exandria & ex-ANN-dree-uh & The world of Critical Role\\
Gaeya Iliera & GUY-yaw ill-ee-AIR-uh & Drow chef in the Aurora Watch\\
Galsariad Ardyth & gal-SAIR-ee-add AR-dith & Drow mage and a member of the rival party\\
Hadarai & HAD-ah-rye & Ghost of an orc priest of Corellon\\
Irvan Wastewalker & UR-vin & Human scoundrel and a member of the rival party\\
Jamil A'alithiya & juh-MEEL AYE-uh-lee-thee-uh & Human High Curator of the Archive of the Cobalt Soul in Ank'Harel\\
Jigow & JEE-gow (rhymes with cow) & Town in Xhorhas populated mostly by goblins and orcs\\
Kalym Telaarin & kal-EEM tel-AW-rin & Drow warrior in Bazzoxan\\
Khelkur the Gull & KEL-ker & Dwarf coleader of the Consortium of the Vermilion Dream\\
Kryn Dynasty & KREEN & Governing body of Xhorhas\\
Larthul the Wolf & LAR-thool & Human coleader of the Consortium of the Vermilion Dream\\
Luxon & LUCKS-awn & Mysterious god of light and rebirth\\
Marquet & mar-KET & Continent covered in deserts, mountains, and jungles\\
Naevyn Tasithar & NAY-vinn TASS-uh-thar & Drow scout in Bazzoxan\\
Nedosi Anay & nuh-DOE-see AH-nay & Half-orc security captain of the Luck's Run casino\\
Perigee & PAIR-uh-jee & Deva agent of Sehanine corrupted by ruidium\\
Postraeck & POST-rake & Human bandit living in the wastes of Xhorhas\\
Prolix Yusaf & PRO-licks & Tiefling agent of the Allegiance of Allsight\\
Rosohna & roh-SO-nuh & Capital of the Kryn Dynasty\\
Ruidium & roo-IDD-ee-um & Corrupting element born of the Netherdeep\\
Ruidus & ROO-idd-iss & Exandria's smaller, vermilion moon\\
Saqiri & sah-KEE-ree & Drow hunter from Alyxian's past\\
Theo Nathope & THEE-oh NAH-tho-pee & A young, innocent avatar of the Apotheon\\
Ushru & OO-shroo & Orc elder of Jigow\\
Verin Thelyss & VAIR-in THAY-liss & Drow military commander of Bazzoxan\\
Watcher Sylralei & SEEL-ruh-lay & Dwarf berserker of the Sentinels of Memory\\
Wildemount & WILD-mount (long i) & Continent covered in vegetation and snow\\
Xhorhas & ZJOR-hawss & Eastern realm of Wildemount, mostly marshes and grasslands\\
Xot & ZJOT & Goliath archaeologist of the Allegiance of Allsight\\
\end{DndTable}

\clearpage
\tikz[remember picture, overlay] \node[inner sep=0pt] at (current page.center){\includegraphics[height=\paperheight]{images/adventure/CRCotN/008-01-001.intro-splash.png}};
\thispagestyle{empty}
\newpage

\chapter{A Fateful Competition}
\DndDropCapLine{L}{ike many of the heroic legends} that echo throughout the annals of Exandrian history, this tale begins in peace—or what passes for peace in Jigow at festival time. Goblin children run through the streets, screaming with mirth, while orcs, humans, drow, and people of countless other folk clink tankards, hurl hatchets at targets, stuff their faces with meat pies, and plunge into the frigid waters of the Ifolon River—all in the name of friendly competition.

The characters are in Jigow during the city's annual Festival of Merit. Though they are the protagonists of this story, today they are but a small handful of excited festivalgoers amid a merry throng. Many of the Xhorhasians present are famed warriors, skilled athletes, cunning con artists, and talented arcanists. Among those taking part in the festivities are the characters' rivals, another group of adventurers whose fates are entwined with the characters.
At the festival's end, the characters and their rivals are chosen to dive into the waters of the Emerald Gulch in search of a prize. Little do they know that this plunge will start them on an adventure greater than any of them could possibly imagine.
\section{Running This Chapter}
This chapter has two parts: the Festival of Merit and the Emerald Grotto. Each part can be played in a single 3- to 4-hour session. Be sure to review the information about the rival nonplayer characters in the "Rivals" section of the introduction before running this chapter. The rivals' interactions with the characters in this chapter will lay the groundwork for relationships that could last throughout the adventure.
The first half of the chapter is a lighthearted and free-form exploration of Jigow and the Festival of Merit. Familiarize yourself with the map of Jigow and the locations keyed to it so that you can effectively guide the characters as they explore the settlement and partake in the festival games.
The second part is a dangerous but straightforward dungeon delve into the Emerald Grotto near Jigow. The characters must overcome a series of challenges to reach the treasure at the end of the dungeon before the rivals do.
\subsection{Rival Impressions}
Whenever a rival is present for an event in this chapter, a "Rival Impressions" section explains how that individual might react to the characters' actions. These suggestions will help you develop the relationships between the rivals and the characters. The evolution of a fierce rivalry into a friendship—or even a romance—is an exciting way to enhance the ongoing story of this adventure. Conversely, a friendship crumbling because of a breakdown in communication or a clash of two unbending ideals is a tragedy that echoes the darker tones of this adventure.
A character can try to change a rival's attitude through roleplaying, which you can adjudicate based on the social interaction rules in the Dungeon Master's Guide. You can also improvise and adjust the rivals' attitudes at your discretion, based on how you feel the rivals would react in any given situation.
\section{Jigow}
Jigow is located in northern Xhorhas on the banks of the Emerald Gulch and the Ifolon River. An amalgamation of coastal villages that were originally settled by several nomadic clans of orcs and goblins, Jigow is now home to humans and other folk as well, and the city gained a significant drow population after Jigow became a part of the Kryn Dynasty.
The villages and townships that make up Jigow are loosely divided into three major areas: the Meatwaters, the main dock area on the shores of the Ifolon River; the Wetwalks, a collection of houses on stilts closest to the wetlands and marshes; and the Jumble, the most densely populated region of the city, where houses are built among giant mangrove trees or on the backs of horizonback tortoises and used as traveling homes.
Jigow has a council of elders who help its people resolve internal conflicts and a military presence in the form of the Aurora Watch—the soldiers of the Kryn Dynasty. The dynasty also has a political liaison in Jigow: a drow named Durth Mirimm (see area J9). Serving as an intermediary between Jigow's elders and the court of the Bright Queen Leylas Kryn, Durth cares more about what Jigow can do for the dynasty than how he can improve the lives of Jigow's citizens.
The Jumble is the heart of Jigow and home to a colorful patchwork of local cultures. The traditions of ancient Xhorhasian peoples, many of them orcs and goblins, live on in the peaceful chaos of the Jumble. The twisting, unplanned streets give the district its name—a medley of homes, shops, workshops, athletic fields, amphitheaters, shrines, taverns, stables, and public spaces for lively gatherings.
Jigow was built from the sweat of its founders' brows, the strength of their muscles, and the music of their voices. In Jigow, any chance to show off one's physical, social, or intellectual prowess is a welcome one, and there's no opportunity better than the annual festival. On this day, friendly and not-so-friendly contests between residents are settled. Every public space is filled with festival games, food stalls, or other amenities for the merrymakers, and plenty of shops and private residences have flung their doors wide to join in the revelry.
\section{Adventure Start}
This story begins as the characters enter the Jumble. Read or paraphrase the boxed text below to set the scene:
\begin{DndReadAloud}{}
You have entered the Jumble, a large district of tangled roads and single-story buildings in the town of Jigow. Throngs of people, most of them orcs and goblins, move through the streets, laughing, singing, running, and sightseeing. All are enthralled by the raucous sights and sounds of the town's Festival of Merit.
You hear snippets of conversations as people pass by: a goblin mother telling her children not to go near the baby horizonback tortoises, a drow guard in shining insectile armor complaining to his partner that his gauntlet was crushed by a hulking orc while arm-wrestling, and a pair of young orcs in swimwear hollering as they rush toward the banks of the Ifolon River.
All around you, colorful signs and banners point toward festival booths surrounded by cheering people. On this street alone, you can see a meat-pie eating contest near a shop mounted on the back of a massive tortoise, and on the other side of the road, a banner emblazoned with the words "Riddles and Rhymes: Unbeatable Riddles!" That banner points toward a three-story temple in the center of the Jumble. The town is yours to explore—where do you want to go?
\end{DndReadAloud}

\subsection{Festival of Merit}
Seven festival games and challenges are described in this chapter, as follows:

\begin{description}[nolistsep]
\item[J1: Best Pies in the Jumble.]
The Unbroken Tusk inn boasts the best meat pies east of the Ashkeeper Peaks. Its owner has offered a prize to the person who can eat the most pies.
\item[J2: One-Shot Solution.]
A maze has been constructed in the Jumble. Competitors win a prize if they navigate the maze in one try without coming upon a dead end.
\item[J3: The Ifolon Plunge.]
Some distance offshore in the Ifolon River, a rusted spear is lodged into a weathered piece of a pier that has long since rotted away. It is a popular spot for Jigow's most talented swimmers to test their speed and skill.
\item[J4: Call to Arms.]
One of the Aurora Watch soldiers stationed in Jigow has challenged anyone and everyone to best her in an arm-wrestling contest. So far, she remains undefeated.
\item[J5: Wetwalks Paddywhack.]
Members of the Gakthash and Uvuroh goblin clans time their rice paddy harvests to coincide with the Festival of Merit every year. This event turns the mundane act of harvesting into a spirited competition.
\item[J6: Herding the Horizonbacks.]
This year's brood of young horizonback tortoises must be relocated from the Wetwalks to the Jumble.
\item[J7: Riddles and Rhymes.]
One of Jigow's elders is fond of making up riddles. He has come up with several new ones for the festival this year.
\end{description}

Each of these challenges is described in greater detail in its own section. Except where otherwise noted, any number of characters can participate in an event.
The characters might want to look for a particular type of game: a contest of athleticism, a test of puzzle-solving, or a trial of endurance. In that case, they must spend a half-hour strolling through the streets in search of the type of challenge they're looking for.
\subsubsection{Medals of Merit Cards}
A character who wins a festival contest earns one of seven different magical medals. Each medal is described in appendix B, and its description also appears on a card in appendix C. You can hand a copy of this card to each player whose character won the medal.
\subsection{Locations in Jigow}
The following locations are keyed to the map of Jigow. They include places where Festival of Merit contests are staged (areas J1 through J7) and other important locations that the characters might visit (areas J8 through J10).

\begin{figure*}
\includegraphics[width=\textwidth]{images/adventure/CRCotN/009-map-1.1-Jigow.png}
\end{figure*}


\subsubsection{J1: Best Pies in the Jumble}
This contest tests the participants' perseverance and endurance. Read:
\begin{DndReadAloud}{}
The savory scents of meat, pastry, and spices fill the air around a three-story building mounted on the back of a gigantic tortoise. A cooking stand and a festival stage with a long table are set up at the tortoise's feet, where an orc stands over a massive oven. She bustles from the oven to the stage and back, placing delicious-looking hand pies at each seat and stacking more of them on a cooling rack nearby.
People are already gathered on stage, including a scrawny young human with a mop of brown hair and a scruffy beard. As you approach, the orc calls out in a melodic baritone, "Come to sample the best meat pies this side of the Wastes, yes? I sell 'em for meals up there, but I'm running pie-eating contests down here all day, if you've got an orc-sized stomach."
\end{DndReadAloud}

Agathe Silverspoon is a chaotic good orc who moves from counter to counter, deftly slicing a variety of cooked and spiced meats, then folding them into piecrusts that she places into the oven behind her.
Agathe gestures toward the raised platform next to her cooking station, where three people are already seated: a cocky male drow flexing his muscles, pointing at people in the gathering crowd and winking; a female halfling (use the thug stat block), sitting very still with a smirk on her face; and a young male human, scanning the crowd as if looking for someone. This individual is Irvan Wastewalker (tier 1), whose stat block appears in appendix A. These contestants have the following modifiers to their Constitution checks and saving throws, which come into play during the pie-eating contest:

\begin{description}[nolistsep]
\item[Male Drow.]
+0
\item[Female Halfling.]
+2
\item[Irvan Wastewalker.]
+1 on checks, +3 on saving throws
\end{description}

\paragraph{Putting Away the Pies}
It costs 5 sp to enter this contest. On Agathe's signal, everyone seated at the platform can begin eating. The pies taste great, the meat seasoned to perfection and the crusts light and fluffy. Of course, there's not much chance to savor the flavor of the pies when the goal is to put away as many of them as possible.
All contestants can easily eat their first pie, requiring no check. Eating a second pie requires a successful DC 8 Constitution check. For every pie after the second, the DC increases by 2. The person who eats the most pies before failing a check is the winner (ties are possible). A contestant whose check fails is too full to continue eating and must make a DC 15 Constitution saving throw. On a failed save, the contestant becomes poisoned for 1 hour and feels queasy.

\end{document}
